\section{AHSME 1967}
        \begin{enumerate}
        	\item (\textit{AHSME 1967}) The three-digit number $2a3$ is added to the number $326$ to give the three-digit number $5b9$.  If $5b9$ is divisible by 9, then $a+b$ equals


 $\text{(A)}\ 2\qquad\text{(B)}\ 4\qquad\text{(C)}\ 6\qquad\text{(D)}\ 8\qquad\text{(E)}\ 9$


	\item (\textit{AHSME 1967}) An equivalent of the expression


 $\left(\frac{x^2+1}{x}\right)\left(\frac{y^2+1}{y}\right)+\left(\frac{x^2-1}{y}\right)\left(\frac{y^2-1}{x}\right)$, $xy \not= 0$,


 is:


 $\text{(A)}\ 1\qquad\text{(B)}\ 2xy\qquad\text{(C)}\ 2x^2y^2+2\qquad\text{(D)}\ 2xy+\frac{2}{xy}\qquad\text{(E)}\ \frac{2x}{y}+\frac{2y}{x}$


	\item (\textit{AHSME 1967}) The side of an equilateral triangle is $s$.  A circle is inscribed in the triangle and a square is inscribed in the circle.  The area of the square is:


 $\text{(A)}\ \frac{s^2}{24}\qquad\text{(B)}\ \frac{s^2}{6}\qquad\text{(C)}\ \frac{s^2\sqrt{2}}{6}\qquad\text{(D)}\ \frac{s^2\sqrt{3}}{6}\qquad\text{(E)}\ \frac{s^2}{3}$


	\item (\textit{AHSME 1967}) Given $\frac{\log{a}}{p}=\frac{\log{b}}{q}=\frac{\log{c}}{r}=\log{x}$, all logarithms to the same base and $x \not= 1$.  If $\frac{b^2}{ac}=x^y$, then $y$ is:


 $\text{(A)}\ \frac{q^2}{p+r}\qquad\text{(B)}\ \frac{p+r}{2q}\qquad\text{(C)}\ 2q-p-r\qquad\text{(D)}\ 2q-pr\qquad\text{(E)}\ q^2-pr$


	\item (\textit{AHSME 1967}) A triangle is circumscribed about a circle of radius $r$ inches.  If the perimeter of the triangle is $P$ inches and the area is $K$ square inches, then $\frac{P}{K}$ is:


 $\text{(A)}\text{independent of the value of} \; r\qquad\text{(B)}\ \frac{\sqrt{2}}{r}\qquad\text{(C)}\ \frac{2}{\sqrt{r}}\qquad\text{(D)}\ \frac{2}{r}\qquad\text{(E)}\ \frac{r}{2}$


	\item (\textit{AHSME 1967}) If $f(x)=4^x$ then $f(x+1)-f(x)$ equals:


 $\text{(A)}\ 4\qquad\text{(B)}\ f(x)\qquad\text{(C)}\ 2f(x)\qquad\text{(D)}\ 3f(x)\qquad\text{(E)}\ 4f(x)$


	\item (\textit{AHSME 1967}) If $\frac{a}{b}<-\frac{c}{d}$ where $a$, $b$, $c$, and $d$ are real numbers and $bd \not= 0$, then:


 $\text{(A)}\ a \; \text{must be negative} \qquad \text{(B)}\ a \; \text{must be positive} \qquad$
$\text{(C)}\ a \; \text{must not be zero} \qquad \text{(D)}\ a \; \text{can be negative or zero, but not positive } \\ \text{(E)}\ a \; \text{can be positive, negative, or zero}$


	\item (\textit{AHSME 1967}) To $m$ ounces of a $m\%$ solution of acid, $x$ ounces of water are added to yield a $(m-10)\%$ solution.  If $m>25$, then $x$ is 


 $\textbf{(A)}\ \frac{10m}{m-10} \qquad \textbf{(B)}\ \frac{5m}{m-10} \qquad \textbf{(C)}\ \frac{m}{m-10} \qquad \textbf{(D)}\ \frac{5m}{m-20} \\ \textbf{(E)}\ \text{not determined by the given information}$


	\item (\textit{AHSME 1967}) Let $K$, in square units, be the area of a trapezoid such that the shorter base, the altitude, and the longer base, in that order, are in arithmetic progression.  Then:


 $\textbf{(A)}\ K \; \text{must be an integer} \qquad \textbf{(B)}\ K \; \text{must be a rational fraction} \\ \textbf{(C)}\ K \; \text{must be an irrational number} \qquad \textbf{(D)}\ K\;  \text{must be an integer or a rational fraction} \qquad$
$\textbf{(E)}\ \text{taken alone neither} \; \textbf{(A)} \; \text{nor} \; \textbf{(B)} \; \text{nor} \; \textbf{(C)} \; \text{nor} \; \textbf{(D)} \; \text{is true}$


	\item (\textit{AHSME 1967}) If $\frac{a}{10^x-1}+\frac{b}{10^x+2}=\frac{2 \cdot 10^x+3}{(10^x-1)(10^x+2)}$ is an identity for positive rational values of $x$, then the value of $a-b$ is:


 $\textbf{(A)}\ \frac{4}{3} \qquad \textbf{(B)}\ \frac{5}{3} \qquad \textbf{(C)}\ 2 \qquad \textbf{(D)}\ \frac{11}{4} \qquad \textbf{(E)}\ 3$


	\item (\textit{AHSME 1967}) If the perimeter of rectangle $ABCD$ is $20$ inches, the least value of diagonal $\overline{AC}$, in inches, is:


 $\textbf{(A)}\ 0\qquad \textbf{(B)}\ \sqrt{50}\qquad \textbf{(C)}\ 10\qquad \textbf{(D)}\ \sqrt{200}\qquad \textbf{(E)}\ \text{none of these}$


	\item (\textit{AHSME 1967}) If the (convex) area bounded by the x-axis and the lines $y=mx+4$, $x=1$, and $x=4$ is $7$, then $m$ equals:


 $\textbf{(A)}\ -\frac{1}{2}\qquad \textbf{(B)}\ -\frac{2}{3}\qquad \textbf{(C)}\ -\frac{3}{2} \qquad \textbf{(D)}\ -2 \qquad \textbf{(E)}\ \text{none of these}$


	\item (\textit{AHSME 1967}) A triangle $ABC$ is to be constructed given a side $a$ (opposite angle $A$). angle $B$, and $h_c$, the altitude from $C$.  If $N$ is the number of noncongruent solutions, then $N$
$\textbf{(A)}\ \text{is} \; 1\qquad \textbf{(B)}\ \text{is} \; 2\qquad \textbf{(C)}\ \text{must be zero}\qquad \textbf{(D)}\ \text{must be infinite}\qquad \textbf{(E)}\ \text{must be zero or infinite}$


	\item (\textit{AHSME 1967}) Let $f(t)=\frac{t}{1-t}$, $t \not= 1$.  If $y=f(x)$, then $x$ can be expressed as 


 $\textbf{(A)}\ f\left(\frac{1}{y}\right)\qquad \textbf{(B)}\ -f(y)\qquad \textbf{(C)}\ -f(-y)\qquad \textbf{(D)}\ f(-y)\qquad \textbf{(E)}\ f(y)$


	\item (\textit{AHSME 1967}) The difference in the areas of two similar triangles is $18$ square feet, and the ratio of the larger area to the smaller is the square of an integer.  The area of the smaller triangle, in square feet, is an integer, and one of its sides is $3$ feet.  The corresponding side of the larger triangle, in feet, is:


 $\textbf{(A)}\ 12\quad \textbf{(B)}\ 9\qquad \textbf{(C)}\ 6\sqrt{2}\qquad \textbf{(D)}\ 6\qquad \textbf{(E)}\ 3\sqrt{2}$


	\item (\textit{AHSME 1967}) Let the product $(12)(15)(16)$, each factor written in base $b$, equal $3146$ in base $b$.  Let $s=12+15+16$, each term expressed in base $b$.  Then $s$, in base $b$, is 


 $\textbf{(A)}\ 43\qquad \textbf{(B)}\ 44\qquad \textbf{(C)}\ 45\qquad \textbf{(D)}\ 46\qquad \textbf{(E)}\ 47$


	\item (\textit{AHSME 1967}) If $r_1$ and $r_2$ are the distinct real roots of $x^2+px+8=0$, then it must follow that:


 $\textbf{(A)}\ |r_1+r_2|>4\sqrt{2}\qquad \textbf{(B)}\ |r_1|>3 \; \text{or} \; |r_2| >3 \\ \textbf{(C)}\ |r_1|>2 \; \text{and} \; |r_2|>2\qquad \textbf{(D)}\ r_1<0 \; \text{and} \; r_2<0\qquad \textbf{(E)}\ |r_1+r_2|<4\sqrt{2}$


	\item (\textit{AHSME 1967}) If $x^2-5x+6<0$ and $P=x^2+5x+6$ then


 $\textbf{(A)}\ P \; \text{can take any real value} \qquad \textbf{(B)}\ 20<P<30\\ \textbf{(C)}\ 0<P<20 \qquad \textbf{(D)}\ P<0 \qquad \textbf{(E)}\ P>30$


	\item (\textit{AHSME 1967}) The area of a rectangle remains unchanged when it is made $2 \frac{1}{2}$ inches longer and $\frac{2}{3}$ inch narrower, or when it is made $2 \frac{1}{2}$ inches shorter and $\frac{4}{3}$ inch wider.  Its area, in square inches, is:


 $\textbf{(A)}\ 30\qquad \textbf{(B)}\ \frac{80}{3}\qquad \textbf{(C)}\ 24\qquad \textbf{(D)}\ \frac{45}{2}\qquad \textbf{(E)}\ 20$


	\item (\textit{AHSME 1967}) A circle is inscribed in a square of side $m$, then a square is inscribed in that  circle, then a circle is inscribed in the latter square, and so on.  If $S_n$ is the sum of the areas of the first $n$ circles so inscribed, then, as $n$ grows beyond all bounds, $S_n$ approaches:


 $\textbf{(A)}\ \frac{\pi m^2}{2}\qquad \textbf{(B)}\ \frac{3\pi m^2}{8}\qquad \textbf{(C)}\ \frac{\pi m^2}{3}\qquad \textbf{(D)}\ \frac{\pi m^2}{4}\qquad \textbf{(E)}\ \frac{\pi m^2}{8}$


	\item (\textit{AHSME 1967}) In right triangle $ABC$ the hypotenuse $\overline{AB}=5$ and leg $\overline{AC}=3$.  The bisector of angle $A$ meets the opposite side in $A_1$.  A second right triangle $PQR$ is then constructed with hypotenuse $\overline{PQ}=A_1B$ and leg $\overline{PR}=A_1C$.  If the bisector of angle $P$ meets the opposite side in $P_1$, the length of $PP_1$ is:


 $\textbf{(A)}\ \frac{3\sqrt{6}}{4}\qquad \textbf{(B)}\ \frac{3\sqrt{5}}{4}\qquad \textbf{(C)}\ \frac{3\sqrt{3}}{4}\qquad \textbf{(D)}\ \frac{3\sqrt{2}}{2}\qquad \textbf{(E)}\ \frac{15\sqrt{2}}{16}$


	\item (\textit{AHSME 1967}) For natural numbers, when $P$ is divided by $D$, the quotient is $Q$ and the remainder is $R$.  When $Q$ is divided by $D'$, the quotient is $Q'$ and the remainder is $R'$.  Then, when $P$ is divided by $DD'$, the remainder is:


 $\textbf{(A)}\ R+R'D\qquad \textbf{(B)}\ R'+RD\qquad \textbf{(C)}\ RR'\qquad \textbf{(D)}\ R\qquad \textbf{(E)}\ R'$


	\item (\textit{AHSME 1967}) If $x$ is real and positive and grows beyond all bounds, then $\log_3{(6x-5)}-\log_3{(2x+1)}$ approaches:


 $\textbf{(A)}\ 0\qquad \textbf{(B)}\ 1\qquad \textbf{(C)}\ 3\qquad \textbf{(D)}\ 4\qquad \textbf{(E)}\ \text{no finite number}$


	\item (\textit{AHSME 1967}) The number of solution-pairs in the positive integers of the equation $3x+5y=501$ is:


 $\textbf{(A)}\ 33\qquad \textbf{(B)}\ 34\qquad \textbf{(C)}\ 35\qquad \textbf{(D)}\ 100\qquad \textbf{(E)}\ \text{none of these}$


	\item (\textit{AHSME 1967}) For every odd number $p>1$ we have:


 $\textbf{(A)}\ (p-1)^{\frac{1}{2}(p-1)}-1 \; \text{is divisible by} \; p-2\qquad \textbf{(B)}\ (p-1)^{\frac{1}{2}(p-1)}+1 \; \text{is divisible by} \; p\\ \textbf{(C)}\ (p-1)^{\frac{1}{2}(p-1)} \; \text{is divisible by} \; p\qquad \textbf{(D)}\ (p-1)^{\frac{1}{2}(p-1)}+1 \; \text{is divisible by} \; p+1\\ \textbf{(E)}\ (p-1)^{\frac{1}{2}(p-1)}-1 \; \text{is divisible by} \; p-1$


	\item (\textit{AHSME 1967}) If one uses only the tabular information $10^3=1000$, $10^4=10,000$, $2^{10}=1024$, $2^{11}=2048$, $2^{12}=4096$, $2^{13}=8192$, then the strongest statement one can make for $\log_{10}{2}$ is that it lies between:


 $\textbf{(A)}\ \frac{3}{10} \; \text{and} \; \frac{4}{11}\qquad \textbf{(B)}\ \frac{3}{10} \; \text{and} \; \frac{4}{12}\qquad \textbf{(C)}\ \frac{3}{10} \; \text{and} \; \frac{4}{13}\qquad \textbf{(D)}\ \frac{3}{10} \; \text{and} \; \frac{40}{132}\qquad \textbf{(E)}\ \frac{3}{11} \; \text{and} \; \frac{40}{132}$


	\item (\textit{AHSME 1967}) Two candles of the same length are made of different materials so that one burns out completely at a uniform rate in $3$ hours and the other in $4$ hours.  At what time P.M. should the candles be lighted so that, at 4 P.M., one stub is twice the length of the other?


 $\textbf{(A)}\ 1:24\qquad \textbf{(B)}\ 1:28\qquad \textbf{(C)}\ 1:36\qquad \textbf{(D)}\ 1:40\qquad \textbf{(E)}\ 1:48$


	\item (\textit{AHSME 1967}) Given the two hypotheses: $\text{I}$ Some Mems are not Ens and $\text{II}$ No Ens are Veens.  If "some" means "at least one," we can conclude that:


 $\textbf{(A)}\ \text{Some Mems are not Veens}\qquad \textbf{(B)}\ \text{Some Vees are not Mems}\\ \textbf{(C)}\ \text{No Mem is a Vee}\qquad \textbf{(D)}\ \text{Some Mems are Vees}\\ \textbf{(E)}\ \text{Neither} \; \textbf{(A)} \; \text{nor} \; \textbf{(B)} \; \text{nor} \; \textbf{(C)} \; \text{nor} \; \textbf{(D)} \; \text{is deducible from the given statements}$


	\item (\textit{AHSME 1967}) $\overline{AB}$ is a diameter of a circle.  Tangents $\overline{AD}$ and $\overline{BC}$ are drawn so that $\overline{AC}$ and $\overline{BD}$ intersect in a point on the circle.  If $\overline{AD}=a$ and $\overline{BC}=b$, $a \not= b$, the diameter of the circle is:


 $\textbf{(A)}\ |a-b|\qquad \textbf{(B)}\ \frac{1}{2}(a+b)\qquad \textbf{(C)}\ \sqrt{ab} \qquad \textbf{(D)}\ \frac{ab}{a+b}\qquad \textbf{(E)}\ \frac{1}{2}\frac{ab}{a+b}$


	\item (\textit{AHSME 1967}) A dealer bought $n$ radios for $d$ dollars, $d$ a positive integer.  He contributed two radios to a community bazaar at half their cost.  The rest he sold at a profit of $8 on each radio sold.  If the overall profit was $72, then the least possible value of $n$ for the given information is:


 $\textbf{(A)}\ 18\qquad \textbf{(B)}\ 16\qquad \textbf{(C)}\ 15\qquad \textbf{(D)}\ 12\qquad \textbf{(E)}\ 11$


	\item (\textit{AHSME 1967}) Let $D=a^2+b^2+c^2$, where $a$, $b$, are consecutive integers and $c=ab$.  Then $\sqrt{D}$ is:


 $\textbf{(A)}\ \text{always an even integer}\qquad \textbf{(B)}\ \text{sometimes an odd integer, sometimes not}\\ \textbf{(C)}\ \text{always an odd integer}\qquad \textbf{(D)}\ \text{sometimes rational, sometimes not}\\ \textbf{(E)}\ \text{always irrational}$


	\item (\textit{AHSME 1967}) In quadrilateral $ABCD$ with diagonals $\overline{AC}$ and $\overline{BD}$ intersecting at $O$, $\overline{BO}=4$, $\overline{OD}=6$, $\overline{AO}=8$, $\overline{OC}=3$, and $\overline{AB}=6$.  The length of $\overline{AD}$ is:


 $\textbf{(A)}\ 9\qquad \textbf{(B)}\ 10\qquad \textbf{(C)}\ 6\sqrt{3}\qquad \textbf{(D)}\ 8\sqrt{2}\qquad \textbf{(E)}\ \sqrt{166}$


	\item (\textit{AHSME 1967}) 
                    \begin{figure}[H]
                        \centering
                        \includegraphics[width=0.5\textwidth]{Problem/AHSME_I_1967_33_Problem_Figure_1.png}
                    \end{figure}
                    


 In this diagram semi-circles are constructed on diameters $\overline{AB}$, $\overline{AC}$, and $\overline{CB}$, so that they are mutually tangent.  If $\overline{CD} \bot \overline{AB}$, then the ratio of the shaded area to the area of a circle with $\overline{CD}$ as radius is:


 $\textbf{(A)}\ 1:2\qquad \textbf{(B)}\ 1:3\qquad \textbf{(C)}\ \sqrt{3}:7\qquad \textbf{(D)}\ 1:4\qquad \textbf{(E)}\ \sqrt{2}:6$


	\item (\textit{AHSME 1967}) Points $D$, $E$, $F$ are taken respectively on sides $AB$, $BC$, and $CA$ of triangle $ABC$ so that $AD:DB=BE:CE=CF:FA=1:n$.  The ratio of the area of triangle $DEF$ to that of triangle $ABC$ is:


 $\textbf{(A)}\ \frac{n^2-n+1}{(n+1)^2}\qquad \textbf{(B)}\ \frac{1}{(n+1)^2}\qquad \textbf{(C)}\ \frac{2n^2}{(n+1)^2}\qquad \textbf{(D)}\ \frac{n^2}{(n+1)^2}\qquad \textbf{(E)}\ \frac{n(n-1)}{n+1}$


	\item (\textit{AHSME 1967}) The roots of $64x^3-144x^2+92x-15=0$ are in arithmetic progression.  The difference between the largest and smallest roots is:


 $\textbf{(A)}\ 2\qquad \textbf{(B)}\ 1\qquad \textbf{(C)}\ \frac{1}{2}\qquad \textbf{(D)}\ \frac{3}{8}\qquad \textbf{(E)}\ \frac{1}{4}$


	\item (\textit{AHSME 1967}) Given a geometric progression of five terms, each a positive integer less than $100$.  The sum of the five terms is $211$.  If $S$ is the sum of those terms in the progression which are squares of integers, then $S$ is:


 $\textbf{(A)}\ 0\qquad \textbf{(B)}\ 91\qquad \textbf{(C)}\ 133\qquad \textbf{(D)}\ 195\qquad \textbf{(E)}\ 211$


	\item (\textit{AHSME 1967}) Segments $AD=10$, $BE=6$, $CF=24$ are drawn from the vertices of triangle $ABC$, each perpendicular to a straight line $RS$, not intersecting the triangle.  Points $D$, $E$, $F$ are the intersection points of $RS$ with the perpendiculars.  If $x$ is the length of the perpendicular segment $GH$ drawn to $RS$ from the intersection point $G$ of the medians of the triangle, then $x$ is:


 $\textbf{(A)}\ \frac{40}{3}\qquad \textbf{(B)}\ 16\qquad \textbf{(C)}\ \frac{56}{3}\qquad \textbf{(D)}\ \frac{80}{3}\qquad \textbf{(E)}\ \text{undetermined}$


	\item (\textit{AHSME 1967}) Given a set $S$ consisting of two undefined elements "pib" and "maa", and the four postulates: $P_1$: Every pib is a collection of maas, $P_2$: Any two distinct pibs have one and only one maa in common, $P_3$: Every maa belongs to two and only two pibs, $P_4$: There are exactly four pibs.  Consider the three theorems: $T_1$: There are exactly six maas, $T_2$: There are exactly three maas in each pib, $T_3$: For each maa there is exactly one other maa not in the same pib with it.  The theorems which are deducible from the postulates are:


 $\textbf{(A)}\ T_3 \; \text{only}\qquad \textbf{(B)}\ T_2 \; \text{and} \; T_3 \; \text{only} \qquad \textbf{(C)}\ T_1 \; \text{and} \; T_2 \; \text{only}\\ \textbf{(D)}\ T_1 \; \text{and} \; T_3 \; \text{only}\qquad \textbf{(E)}\ \text{all}$


	\item (\textit{AHSME 1967}) Given the sets of consecutive integers $\{1\}$,$\{2, 3\}$,$\{4,5,6\}$,$\{7,8,9,10\}$,$\; \cdots \;$, where each set contains one more element than the preceding one, and where the first element of each set is one more than the last element of the preceding set.  Let $S_n$ be the sum of the elements in the nth set.  Then $S_{21}$ equals:


 $\textbf{(A)}\ 1113\qquad \textbf{(B)}\ 4641 \qquad \textbf{(C)}\ 5082\qquad \textbf{(D)}\ 53361\qquad \textbf{(E)}\ \text{none of these}$


	\item (\textit{AHSME 1967}) Located inside equilateral triangle $ABC$ is a point $P$ such that $PA=8$, $PB=6$, and $PC=10$.  To the nearest integer the area of triangle $ABC$ is:


 $\textbf{(A)}\ 159\qquad \textbf{(B)}\ 131\qquad \textbf{(C)}\ 95\qquad \textbf{(D)}\ 79\qquad \textbf{(E)}\ 50$


\end{enumerate}